\subsection{Biased couplings of product coordinates}
\label{p:sec:couplings}

The measures $\lambda_I$ give the pair identity when only finitely many
coordinates are replaced.  We now construct ordered probability laws
that give every block opportunities to cross its thresholds.  Their
endpoint densities remain uniformly bounded in $L^2(\mu)$, and the
crossing estimates remain valid after either entire endpoint is fixed.
These are the properties needed for the convolution obstruction.

An ordered coupling with both marginals equal to Haar measure would
give $x=y$ almost surely: in every coordinate,
$y_i-x_i\ge0$ and $\mathbb E(y_i-x_i)=0$.  We therefore allow a controlled
change in the endpoint marginals.  Fix a block index $j\ge1$.  For
$i\in J_s$, write $m=m_s$ and set
\begin{equation}\label{p:eq:rho-definition}
 \rho_s(w)=\kappa_s\sqrt{w/m},\qquad w=a_s,b_s.
\end{equation}
Choose the coordinate pair $(x_i,y_i)$ by the following mixture:
\begin{itemize}
 \item For each of the two labelled widths $w=a_s,b_s$, include a
 top component of probability $\rho_s(w)$.  In this component $x_i$
 and $y_i$ are independent uniforms on $[0,1-w/2)$ and
 $[1-w/2,1)$, respectively.
 \item For each labelled width $w=a_s,b_s$, include a bottom component
 of probability $\rho_s(w)$, with independent uniforms on $[0,w/2)$
 and $[w/2,1)$, respectively.
 \item If $s=j$, include a tagged component of probability $1/m$,
 with independent uniforms on $[0,1/m)$ and $[1-1/m,1)$,
 respectively.
 \item With the remaining probability
 \begin{equation}\label{p:eq:diagonal-probability}
  q_{j,s}=1-2\rho_s(a_s)-2\rho_s(b_s)
                  -\ind_{\{s=j\}}/m,
 \end{equation}
 choose $x_i=y_i$ uniformly in $[0,1)$.
\end{itemize}
The width labels remain distinct even if their numerical values agree.
The scale bounds imply
\begin{equation}\label{p:eq:width-bounds}
 \frac1m\le a_s,b_s\le\frac14,
 \qquad
 \frac{\rho_s(w)}{w}=\frac{\kappa_s}{\sqrt{mw}}
       \le\kappa_s\le\frac1{100}.
\end{equation}
Since $m\ge100$, it follows that
$q_{j,s}\ge1-\kappa_s-1/m\ge0.98$.

Make the choices independently over all coordinates.  More precisely,
let $(\Omega_j,\mathbb P_j)$ be the countable product of the spaces
carrying the mixture label and its coordinate pair.  Let
$x,y:\Omega_j\to X$ be the endpoint maps, and let $\nu_j$ be the law of
$(x,y)$ on $X^2$.  The labels, including the tags, are thus retained as
events on $\Omega_j$.  These are the actual product coordinates of the
endpoints.  Each component is supported on $x_i\le y_i$, so $x\le y$
almost surely.

\begin{lemma}[Marginals and conditional crossings]
\label{p:lem:coupling-marginals}
Put
\begin{equation}\label{p:eq:uniform-product-density}
 H^2=\exp\left(20\left(2\sum_{s\ge1}\kappa_s^2+1\right)\right)<\infty.
\end{equation}
For every $j\ge1$, the law just constructed has the following properties.
\begin{enumerate}
 \item Every one-coordinate endpoint density is at most $10$.  Both
 marginals of $\nu_j$ have $L^2(\mu)$ densities of norm at most $H$.
 The same bound holds for the product law formed by any deterministic
 choice of the $x$-marginal or the $y$-marginal at each coordinate.
 \item Conditional on $x$, the coordinates $y_i$ are independent.
 Conditional on $y$, the coordinates $x_i$ are independent.  If
 $i\in J_s$, $w=a_s$ or $b_s$, and $u\in[w,1-w]$, then
 \begin{align}
  \mathbb P_j(x_i<u\le y_i\mid x)&\ge\rho_s(w)/10
       &&\text{on }\{x_i<u\},\label{p:eq:crossing-given-x}\\
  \mathbb P_j(x_i<u\le y_i\mid y)&\ge\rho_s(w)/10
       &&\text{on }\{y_i\ge u\}.
       \label{p:eq:crossing-given-y}
 \end{align}
 The assertions hold almost surely, simultaneously for any specified
 countable family of such thresholds.
\end{enumerate}
\end{lemma}

\begin{proof}
For a top or bottom component of width $w$, each endpoint interval has
length at least $w/2$, so its density is at most $2/w$.  A tagged
endpoint density is at most $m$.  Consequently either one-coordinate
marginal density $h_i$ satisfies
\[
 0<q_{j,s}\le h_i
 \le1+4\sum_{w=a_s,b_s}\frac{\rho_s(w)}w+\ind_{\{s=j\}}
 \le2+8\kappa_s<10.
\]
Throughout these sums the two widths are counted with their labels.

If $v$ is the uniform density on an interval of length at least $w/2$,
then $\int(v-1)=0$ and $\|v-1\|_2^2\le2/w$.  The corresponding
bound for a tagged interval is $m$.  Subtracting $1$ from $h_i$
expresses it as the sum of four weighted mean-zero perturbations and,
when $s=j$, one tagged perturbation.  Cauchy--Schwarz for these at most
five summands gives
\begin{equation}\label{p:eq:marginal-energy}
 \begin{split}
 \|h_i-1\|_2^2
 &\le20\sum_{w=a_s,b_s}\frac{\rho_s(w)^2}{w}
                    +5\frac{\ind_{\{s=j\}}}{m_s}\\
 &\le20\left(\frac{2\kappa_s^2}{m_s}
                    +\frac{\ind_{\{s=j\}}}{m_s}\right).
 \end{split}
\end{equation}
This estimate holds for either endpoint.  Summing over the $m_s$
coordinates of every block gives the bound
$20(2\sum_s\kappa_s^2+1)$ for every deterministic selection of endpoint
marginals.  Each selected coordinate density is nonnegative and has
integral one, and the coordinate laws are independent.  Thus
Lemma~\ref{lem:product-density-bound} supplies the asserted $L^2$ density
with norm at most $H$.  In particular, for each such product law $\eta$
and every Borel set $A\subseteq X$,
\begin{equation}\label{p:eq:uniform-absolute-continuity}
 \eta(A)\le H\mu(A)^{1/2}.
\end{equation}

To verify the assertions after conditioning on an entire endpoint,
we describe the conditional kernels.  At coordinate $i\in J_s$, let
$\ell$ run through the replacement labels, including the tag when
present.  Write $\pi_\ell$ for its mixture probability and
$u_\ell,v_\ell$ for the normalized interval densities of its two
endpoints.  The marginal densities in this coordinate are
\[
 h_i^x(t)=q_{j,s}+\sum_\ell\pi_\ell u_\ell(t),
 \qquad
 h_i^y(t)=q_{j,s}+\sum_\ell\pi_\ell v_\ell(t).
\]
They are strictly positive.  A conditional law of $y_i$ given $x_i=t$
is
\begin{equation}\label{p:eq:conditional-kernel}
 K_i^x(t,dz)=\frac{q_{j,s}}{h_i^x(t)}\delta_t(dz)
       +\sum_\ell\frac{\pi_\ell u_\ell(t)}{h_i^x(t)}
                              v_\ell(z)\,dz.
\end{equation}
The reverse kernel $K_i^y$ interchanges $u_\ell$ and $v_\ell$.
Independence of the coordinate pairs shows, first for cylinder tests,
that $\bigotimes_i K_i^x(x_i,\cdot)$ is a conditional law of $y$ given
$x$.  The monotone class theorem extends the defining identity to
all Borel tests on $X^2$.  The corresponding product of reverse
kernels is a conditional law of $x$ given $y$.  These formulas prove
both assertions of conditional independence.

If $x_i=t<u\le1-w$, then $t\in[0,1-w/2)$, where the top component of
width $w$ has $x_i$-density $(1-w/2)^{-1}\ge1$.  Its whole $y_i$
interval lies above $u$.  Since $h_i^x\le10$, its contribution to
\eqref{p:eq:conditional-kernel} is at least $\rho_s(w)/10$.
This proves \eqref{p:eq:crossing-given-x}.  If $y_i=t\ge u\ge w$,
the bottom component has $y_i$-density $(1-w/2)^{-1}\ge1$ at $t$,
and its whole $x_i$ interval lies below $u$.  The reverse kernel
therefore proves \eqref{p:eq:crossing-given-y}.  Intersecting the
full-measure sets for any specified countable family of thresholds
gives the final assertion.
\end{proof}

It remains to establish the pair identity under these probability laws.
We first replace a finite prefix of $x$ by coordinates of $y$.  The
resulting law is dominated by a finite sum of the triangular laws
$\lambda_I$, but its diagonal tail retains the biased $x$-distribution.
Keeping that tail density explicit permits the transfer to infinitely
many replacements.

\begin{proposition}[The identity under the biased couplings]
\label{p:prop:biased-identity}
Let $p\in M$ be an orthogonal projection satisfying $(1-p)Tp=0$ in the
product model.  For every $j\ge1$,
\begin{equation}\label{p:eq:biased-identity}
 (1-p(y))\operatorname{diag}_{k\in\Z}\bigl(R_{x,y}(k)\bigr)p(x)=0
 \qquad\text{for }\nu_j\text{-almost every }(x,y).
\end{equation}
Moreover, $0<R_{x,y}(k)\le1$ for all $k\in\Z$ almost surely.
\end{proposition}

\begin{proof}
Fix $j$.  For $D\ge1$, define
\[
 y_i^{(D)}=
 \begin{cases}
  y_i,&i\le D,\\
  x_i,&i>D,
 \end{cases}
\]
and let $\beta_D$ be the law of $(x,y^{(D)})$ under $\mathbb P_j$.
We claim that
\begin{equation}\label{p:eq:finite-prefix-domination}
 \beta_D\ll\sum_{I\subseteq\{1,\ldots,D\}}\lambda_I.
\end{equation}
Condition on a pattern $c$ of mixture labels in the first $D$
coordinates.  There are finitely many patterns.  Let $I(c)$ consist
of the coordinates carrying replacement labels.  At each $i\in I(c)$
the conditioned pair has density
\[
 r_{c,i}(t,z)=
 \frac{\ind_{U_{c,i}}(t)\ind_{V_{c,i}}(z)}
 {|U_{c,i}|\,|V_{c,i}|}
\]
relative to the unnormalized triangular measure
$\ind_{\{0\le t<z<1\}}\,dt\,dz$.  Here $U_{c,i},V_{c,i}$ are the
two intervals of the selected component.  Their lengths are positive
and their rectangle is supported on $t<z$.  The remaining prefix
coordinates are uniformly diagonal.

For the diagonal tail, let $g$ be the full $x$-marginal density from
Lemma~\ref{p:lem:coupling-marginals}.  Write
$\mu=\mu_{\le D}\otimes\mu_{>D}$ and define
\begin{equation}\label{p:eq:tail-density}
 g_D(z)=\int g(u,z)\,d\mu_{\le D}(u),
 \qquad
 \|g_D\|_{L^2(\mu_{>D})}\le\|g\|_{L^2(\mu)}\le H.
\end{equation}
Fubini identifies $g_D$ as the density of $x_{>D}$, and Jensen's
inequality proves the norm bound.  Independence across coordinates
ensures that conditioning on $c$ leaves this tail law unchanged.
Therefore the conditional law of $(x,y^{(D)})$ has density
\begin{equation}\label{p:eq:conditioned-prefix-density}
 g_D(x_{>D})\prod_{i\in I(c)}r_{c,i}(x_i,y_i)
\end{equation}
with respect to $\lambda_{I(c)}$.  This follows on cylinder sets
from the coordinate construction, and hence on all Borel sets.
The rectangle factors form a finite product and are bounded.  The
tail factor is integrable and need not equal one.  Summing the
conditional laws over the finitely many patterns proves
\eqref{p:eq:finite-prefix-domination}.

Proposition~\ref{p:prop:pair-identity} now gives the pair identity
$\beta_D$-almost everywhere.  To pass to the limit, the first marginal
of $\beta_D$ is the $x$-marginal of $\nu_j$, while its second marginal
uses the $y$-marginals in the first $D$ coordinates and the
$x$-marginals thereafter.  Lemma~\ref{p:lem:coupling-marginals} and
\eqref{p:eq:uniform-absolute-continuity} therefore give uniform
absolute continuity of both marginal families, and absolute continuity
of the two marginals of $\nu_j$.

On the common probability space, $y^{(D)}\to y$ in the product topology.
Since $X^2$ is compact metrizable, bounded convergence against continuous
functions proves $\beta_D\Rightarrow\nu_j$.  Every one of these laws
is carried by ordered pairs.  Their absolutely continuous marginals
put both endpoints in the common conull set on which the cocycles are
finite and positive.  The ordered-coordinate formula for the cocycle
ratios consequently gives $0<R_{x,y^{(D)}}(k)\le1$ and
$0<R_{x,y}(k)\le1$, simultaneously for every signed $k$.
All hypotheses of Lemma~\ref{lem:compact-transfer} are now satisfied,
with the fixed cap $B=1$.  That lemma transfers the identity from
$\beta_D$ to $\nu_j$, proving \eqref{p:eq:biased-identity}.
\end{proof}
